\section{A buffered horizon and bending on every physical record}
\label{sec:cb-bending}
The physical horizon established in Section~\ref{sec:eh-horizon}
contains a strict covering margin which has not yet been used.  This
margin gives a deterministic lower bound for accumulated reflection,
not just an upper bound for a single flight.  All centers in this
section belong to the original triangular lattice; all disks still
have radius $R\in[R_-,R_+]=[9/20,47/100]$.  Incidence means
$c_i=\sqrt{1-p_i^2}$ at an actual collision.  In particular this
section imposes no new cutoff on the exact event
\eqref{eq:v76-event}.

\begin{lemma}[The unused normal margin in the lattice covering]
\label{lem:cb-buffer}
Set
\begin{equation}\label{eq:cb-constants}
 H=\frac{28}{5},\qquad \omega_0=\frac{179}{400},\qquad
 b_0=R_--\omega_0=\frac1{400},\qquad
 \gamma_0=\frac{b_0}{2H}=\frac1{4480}.
\end{equation}
For every $q\in\R^2$ and every unit vector $v$, there is a center
$z\in\Lambda$ whose perpendicular projection onto
$\{q+tv:0\leq t\leq H\}$ is at a parameter $s\in(0,H)$ and
satisfies $|z-(q+sv)|<\omega_0$.  The assertion is simultaneous
in the starting point and direction, including limiting directions.
\end{lemma}
\begin{proof}
Apply Lemma~\ref{lem:eh-rectangle} at the midpoint $q+Hv/2$.
The four direction intervals and primitive row vectors in the proof
of Theorem~\ref{thm:eh-horizon} cover every direction by lattice
symmetry.  Their squared normal widths are bounded, respectively, by
\[
 \frac{961}{4804},\quad\frac{13697401}{73866304},\quad
 \frac{529}{3136},\quad\frac{75}{448}.
\]
Each is at most the first number, and
\[
 \frac{961}{4804}<\frac{32041}{160000}
                  =\left(\frac{179}{400}\right)^2.
\]
These are comparisons of positive rational numbers.  The absolute
longitudinal displacement from the midpoint is at most
$\tfrac12\sqrt{3847/124}<14/5=H/2$, as proved there.  Thus its foot
lies strictly inside the stated segment.  No simulated orbit or
unstated compactness constant is involved.  The number $\omega_0$
is a covering distance, not a changed scatterer radius.
\end{proof}

\begin{theorem}[Every physical time block carries positive bending]
\label{thm:cb-time-bending}
Let $\gamma:[0,H]\to\R^2$ be any unit-speed trajectory segment of
the original Lorentz gas, with the actual specular reflections at
times $0<t_i<H$.  Then
\begin{equation}\label{eq:cb-time-bending}
 \sum_{0<t_i<H}c_i\geq\gamma_0.
\end{equation}
One may start or end inside a flight.  For a collision at the initial
time, the initial direction is the outgoing one.  The same statement
holds for a fixed finite-word grazing limit which is still a
nonpenetrating specular polygon; a tangential contact has $c_i=0$.
The bound is deterministic and uniform in $R$, the word, the starting
point and the initial direction.
\end{theorem}
\begin{proof}
There are finitely many collisions in a bounded physical time
interval: consecutive flights have length at least
$\tau_-=3/50$, with only the first and last partial flights possibly
shorter.  Write $v_0$ for the initial outgoing velocity and
$\Delta v_i=v(t_i+)-v(t_i-)$.  Specular reflection at a circular
disk with outward unit normal $n_i$ gives
\[
 \Delta v_i=-2\bigl(v(t_i-)\cdot n_i\bigr)n_i,
 \qquad |\Delta v_i|=2c_i.
\]
Integrating the piecewise constant velocity, for $0\leq t\leq H$
we have the exact identity and bound
\begin{equation}\label{eq:cb-polygon-deviation}
 \gamma(t)-\gamma(0)-tv_0
   =\sum_{0<t_i<t}(t-t_i)\Delta v_i,
 \qquad
 |\gamma(t)-\gamma(0)-tv_0|
   \leq 2H\sum_{0<t_i<H}c_i.
\end{equation}
Choose $z$ and $s$ from Lemma~\ref{lem:cb-buffer}.  If the sum in
\eqref{eq:cb-time-bending} were less than $\gamma_0$, then
\[
 |\gamma(s)-z|
 <\omega_0+2H\gamma_0=R_-\leq R.
\]
This places a true physical point inside an open scatterer, a
contradiction.  The argument uses neither disjoint reflection events
nor statistical independence.  It also applies directly to the stated
limiting polygons, for which the same velocity identity holds and
which remain outside every open disk.  It does not define a billiard
continuation through an arbitrary singular endpoint.
\end{proof}

\begin{corollary}[A uniform block bound over all collision words]
\label{cor:cb-collision-bending}
For any actual regular record with $m\geq1$ flights, put $N_0=94$.
Then
\begin{align}
 \sum_{i=1}^{m-1}c_i
   &\geq\gamma_0\left\lfloor\frac{m}{N_0}\right\rfloor,
       \label{eq:cb-all-collisions}\\
 \#\left\{1\leq i<m:c_i\geq\frac1{416640}\right\}
   &\geq\left\lfloor\frac{m}{N_0}\right\rfloor.
       \label{eq:cb-positive-fraction}
\end{align}
Every full block of $94$ consecutive flights supplies its own
collision in the second inequality; these supplied collisions are
distinct.  All other collisions remain in the record and in every
observable.
\end{corollary}
\begin{proof}
The duration of a $94$-flight block is at least
$94\tau_-=141/25>28/5=H$.  Starting at its first collision with the
outgoing velocity, Theorem~\ref{thm:cb-time-bending} supplies a sum
at least $\gamma_0$ among its $93$ interior collision indices.  The
sets of such interior indices for disjoint successive blocks are
disjoint.  Summation proves \eqref{eq:cb-all-collisions}.  Within each
block at least one of the $93$ nonnegative summands is at least
$\gamma_0/93=1/416640$, proving \eqref{eq:cb-positive-fraction}.
The possibly incomplete last block is not assigned any fictitious
contribution.  For $m<94$ both statements retain their stated
meaning, with a zero right side.
\end{proof}

\begin{remark}[What positive bending does and does not measure]
\label{rem:cb-scope}
The result concerns visited incidences along every original physical
record.  It is not abundance of the original comparison pool:
$\mathsf U^{\rm low}_{B,m}$ and $P^{\rm ng}_{r_B,m}$ carry additional
roof, clearance, pairing and exact-label requirements.  A positive
fraction of moderately transverse collisions does not pay a single
unvisited tangency or force any endpoint trace to vanish.  We use the
bound below only in an exact variation of actual critical physical
lengths; we do not multiply an unscaled mass bound by $m^2$ and call
it small.
\end{remark}
